%%=============================================================================
%% Conclusie
%%=============================================================================

\chapter{Conclusie}%
\label{ch:conclusie}

% TODO: Trek een duidelijke conclusie, in de vorm van een antwoord op de
% onderzoeksvra(a)g(en). Wat was jouw bijdrage aan het onderzoeksdomein en
% hoe biedt dit meerwaarde aan het vakgebied/doelgroep? 
% Reflecteer kritisch over het resultaat. In Engelse teksten wordt deze sectie
% ``Discussion'' genoemd. Had je deze uitkomst verwacht? Zijn er zaken die nog
% niet duidelijk zijn?
% Heeft het onderzoek geleid tot nieuwe vragen die uitnodigen tot verder 
%onderzoek?

De centrale onderzoeksvraag van deze bachelorproef luidde: \textit{Hoe kan een datagedreven scoutingomgeving voor de BNXT League worden ontworpen en gerealiseerd, zodat scouts spelers efficiënt kunnen analyseren en vergelijken op basis van relevante prestatie-indicatoren?} Op basis van de literatuurstudie en de realisatie van het proof-of-concept blijkt dat een combinatie van een relationeel datawarehouse in sterschema, een robuuste ETL-pipeline en een interactief Power BI-dashboard een passend antwoord biedt binnen de vooropgestelde scope.

Er werd een volledige scoutingomgeving opgebouwd die bestaat uit een centrale datalaag in Microsoft SQL Server, een ETL-pipeline in Python en een Power BI dashboard met meerdere functionele pagina's. De omgeving laat scouts en data-analisten toe om spelers te filteren op onder meer seizoen, team, positie, leeftijd, lengte en nationaliteit, individuele spelersprofielen te raadplegen met kernstatistieken en een vergelijking met het leaguegemiddelde, spelers rechtstreeks naast elkaar te plaatsen via een vergelijkingspagina en een shortlist van kandidaten op te bouwen via een uitgebreide tabel met filters. Daarmee wordt de kloof verkleind tussen de gefragmenteerde, intuïtieve werkwijze met losse Excel bestanden en websites en een meer gestructureerde, herhaalbare besluitvorming op basis van een centrale datalaag.

\section*{Antwoord op de deelvragen}

De eerste deelvraag betrof de statistieken en prestatie-indicatoren die relevant zijn voor BNXT-scouting. Via de literatuurstudie naar talentidentificatie en key performance indicators in basketbal werden zowel antropometrische kenmerken (zoals lengte, gewicht en leeftijd) als wedstrijdstatistieken (punten, rebounds, assists, steals, blocks, turnovers, shootingpercentages en efficiëntiescores) geselecteerd. Deze indicatoren zijn in het datamodel vertaald naar velden en berekende kolommen, zodat ze rechtstreeks gebruikt kunnen worden in het dashboard voor filtering, rangschikking en vergelijking van spelers.

De tweede deelvraag onderzocht hoe historische data uit verschillende bronnen verzameld, opgeschoond en geïntegreerd kan worden in een centrale datalaag. In dit onderzoek werden publieke bronnen zoals de officiële BNXT League-website en RealGM gecombineerd, waarbij Python-scripts gebruikt werden om zowel spelersprofielen als wedstrijdstatistieken per speler per wedstrijd op te halen. Door systematische normalisatie, fuzzy matching van spelernamen en expliciete mappingtabellen voor team -en spelersnamen konden inconsistenties tussen bronnen weggewerkt worden en werden meer dan tweehonderd naamkoppelingen opgebouwd. Dit toont aan dat het opbouwen van een betrouwbare datalaag in de praktijk een aanzienlijke tijdsinvestering in datakwaliteit vergt, maar dat een consistente set spelers- en teamidentiteiten haalbaar is op basis van publiek beschikbare data.

De derde deelvraag richtte zich op de inzet van machine-learningmodellen voor het rangschikken en vergelijken van spelers. Uit de literatuur blijkt dat onder meer clusteringmethoden geschikt zijn om spelers op basis van speelstijl en statistisch profiel te groeperen en functionele vervangers te identificeren. In deze bachelorproef werd deze benadering theoretisch uitgewerkt, maar niet volledig geïntegreerd in het finale prototype: de rangschikking op de Scouting Shortlist-pagina is gebaseerd op direct sorteerbare statistieken en samengestelde efficiëntiematen. Dit is een bewuste afbakening in functie van de proof-of-concept, en vormt tegelijk een duidelijke groeipiste voor toekomstige versies waarin geavanceerdere modellen kunnen worden toegevoegd.

De vierde deelvraag vroeg hoe een gebruiksvriendelijk Power BI-dashboard ontworpen kan worden om scouts te ondersteunen. Op basis van de requirementsanalyse en richtlijnen uit de literatuur over dashboardontwerp werd een prototype opgebouwd met onder meer een League Overview, een Team Dashboard, een Player Profile, een Player Comparison-pagina en een Scouting Shortlist. Een eerste gebruiksevaluatie met vooraf gedefinieerde scenario's, zoals het identificeren van spelers met een bepaald profiel of het zoeken van een vervanger voor een vertrekkende speler, toonde aan dat deze taken effectief met het dashboard kunnen worden uitgevoerd. De pagina-indeling en filtermogelijkheden sluiten aan bij de informatiebehoeften van scouts en data-analisten zoals beschreven in de literatuur.

De vijfde en laatste deelvraag had betrekking op de meerwaarde van de ontwikkelde omgeving ten opzichte van de huidige, meer intuïtieve benadering van spelersanalyse binnen de BNXT League. Waar in de huidige praktijk vaak met losse Excel-bestanden, verschillende websites en manuele berekeningen wordt gewerkt, biedt de gerealiseerde omgeving één centraal platform met opgeschoonde en gestandaardiseerde data. Dit maakt het eenvoudiger om spelers op een consistente manier te vergelijken, analyses te herhalen en beslissingen transparant te onderbouwen binnen een club. Tegelijk blijft er ruimte voor kwalitatieve beoordeling: de omgeving is opgezet als beslissingsondersteunend hulpmiddel dat de expertise van scouts en coaches aanvult in plaats van vervangt.

\section*{Kritische reflectie}

Hoewel het resultaat grotendeels aansluit bij de verwachtingen, zijn er enkele beperkingen die in een kritische reflectie benoemd moeten worden. Ten eerste werd de machine-learningcomponent niet geïntegreerd in het finale prototype. De literatuur toont aan dat geavanceerde analytische technieken, zoals clustering en modelgebaseerde rangschikkingen, bijkomende inzichten kunnen opleveren in speelstijl en spelersprofielen, maar de realisatie van een correcte datalaag en een volledig functioneel dashboard vergde in deze bachelorproef de meeste aandacht. De keuze om te focussen op een robuust datamodel en een bruikbaar dashboard is te verdedigen vanuit het doel van een proof-of-concept, maar maakt de omgeving op dit moment minder onderscheidend ten opzichte van generieke rapportagetools.

Ten tweede is de gebruiksevaluatie beperkt gebleven tot interne tests aan de hand van vooraf gedefinieerde use cases en een beperkte groep testgebruikers. Een grootschalige gebruikersstudie met scouts, coaches en data-analisten uit verschillende BNXT-clubs zou een objectievere beoordeling opleveren van bruikbaarheid, leercurve en toegevoegde waarde in de praktijk. De literatuur benadrukt dat datagedreven tools pas volledig tot hun recht komen wanneer ze nauw aansluiten bij de bestaande werkprocessen en wanneer eindgebruikers actief betrokken worden bij ontwerp en iteratieve verbetering.

Ten derde beperkt de huidige omgeving zich tot traditionele boxscore-statistieken en een aantal afgeleide efficiëntiematen. Meer geavanceerde indicatoren, zoals usage rate, plus/minus-varianten of contextspecifieke metrics, konden niet volledig worden berekend op basis van de beschikbare ruwe data, bijvoorbeeld door het ontbreken van gedetailleerde play-by-play of trackinginformatie. Dit verkleint de analytische diepgang, zeker bij het vergelijken van spelers uit teams met sterk verschillende tempo's of speelstijlen, en benadrukt de nood aan rijkere datasets voor toekomstig onderzoek.

\section*{Bijdrage aan het onderzoeksdomein}

De bijdrage van deze bachelorproef situeert zich op het snijvlak van data-engineering, business intelligence en sportanalyse. Bestaande literatuur en praktijkcases rond datagedreven scouting focussen vaak op grote competities zoals de NBA of blijven op een abstract procesniveau, waardoor de vertaling naar regionale competities met beperktere middelen minder duidelijk is. Deze bachelorproef toont aan dat een geïntegreerde scoutingomgeving ook haalbaar is voor een competitie zoals de BNXT League, op basis van publiek beschikbare data en courante technologieën zoals Python, Microsoft SQL Server en Power BI.

Het uitgewerkte sterschema, de ETL-pipeline en het dashboardontwerp kunnen dienen als referentiemodel voor andere competities of clubs die een gelijkaardig platform willen opzetten. Daarnaast vormt het proof-of-concept een concreet vertrekpunt voor vervolgonderzoek naar rijkere datasets, geavanceerdere analysetechnieken en nauwere integratie met bestaande clubprocessen. Op die manier levert deze bachelorproef niet alleen een werkbare oplossing voor de BNXT-context, maar ook een bouwsteen voor de verdere professionalisering van datagedreven scouting in regionale basketbalcompetities.